\section{Experimental Setup}

\subsection{Datasets and Directed Transfer Settings}

The reset table system evaluates three Amazon target-domain settings. For
AM-CDs, AM-Electronics and AM-Movies are used as auxiliary domains. For
AM-Electronics, AM-CDs and AM-Movies are used as auxiliary domains. For
AM-Movies, AM-CDs and AM-Electronics are used as auxiliary domains. In each
setting, the target dataset is the evaluation dataset, and the auxiliary domain
is the source of cross-domain transfer signal.

\subsection{Metrics}

We report automatic text-generation metrics and rating-error metrics in the
same overall performance tables. The text metrics are ROUGE-1, ROUGE-L, BLEU-1,
BLEU-2, BLEU-3, BLEU-4, METEOR, DIST-1, and DIST-2. Rating errors are reported
by RMSE and MAE. Larger values are better for the automatic text metrics, while
smaller values are better for RMSE and MAE.

\subsection{Baselines and Comparison Scope}

The main tables follow the D4C-style comparison format. Single-domain
baselines are DualPC, PETER, ReXPlug, and ERRA. Cross-domain baselines are
PEFT, PEPLER, AdaReX, and D4C. Our method is reported as ReCEG in the result
tables. D4C is the closest causal reference point for the dynamic
textual-feedback view \cite{d4c}. Classical explainable recommendation systems
are cited as task-family context where appropriate \cite{peter,pepler,adarex}.

\subsection{Evaluation Scope}

We report current results on three Amazon target-domain settings. The baseline
rows follow D4C-style reporting, but standard deviations are omitted for
compactness. The results should be read as bounded evidence rather than as a
multi-seed benchmark claim.
